% ===============================
% Worksheet / Manual Template
% ===============================
\documentclass[12pt, a4paper]{article}

% ===============================
% Metadata
% ===============================
\newcommand*{\CourseCode}{MAT221}
\newcommand*{\CourseTitle}{Real Analysis}
\newcommand*{\Chapter}{Absolute and Conditional Convergence}
\newcommand*{\Topics}{Rearrangements, Double Summations and Products}
\newcommand*{\DocDate}{December 7, 2025}

\include{header}

% ==========================================
% Document
% ==========================================
\begin{document}
\FrontPage
\Large


\begin{heading}
	Convergence of Infinite Series
\end{heading}

\begin{definition}[Convergence of Infinite Series]
	We define the corresponding sequence of partial sums $(s_m)$ by
	\[
	s_m = a_1 + a_2 + a_3 + \cdots + a_m,
	\]
	and say that the series 
	\(
	\sum_{n=1}^{\infty} a_n
	\)
	converges to $S$ if the sequence $(s_m)$ converges to $S$.  
	In this case, we write
	\[
	\sum_{n=1}^{\infty} a_n = S.
	\]
	\vspace{5pt}
	
	If the sequence $\{S_N\}$ does not converge, the series is called divergent.
\end{definition}

\clearpage

\begin{heading}
	Absolute Convergence
\end{heading}

\vspace{10pt}

\begin{definition}[Absolute Convergence]
	A series $\sum_{n=1}^{\infty} a_n$ is said to be \textit{absolutely convergent} if the series of absolute values $\sum_{n=1}^{\infty} |a_n|$ converges.	
\end{definition}

\vspace{10pt}

\begin{theorem}
	Every absolutely convergent series is convergent.
\end{theorem}

\clearpage

\begin{heading}
	Conditional Convergence
\end{heading}

\vspace{10pt}

\begin{definition}[Conditional Convergence]
	A series $\sum_{n=1}^{\infty} a_n$ is said to be \textit{conditionally convergent} if it converges but does not converge absolutely.
\end{definition}

\clearpage

\begin{problem}
	The alternating harmonic series
	\[
	\sum_{n=1}^{\infty} \frac{(-1)^{n+1}}{n} = 1 - \frac{1}{2} + \frac{1}{3} - \frac{1}{4} + \cdots
	\]
	is conditionally convergent.
\end{problem}

\clearpage

\begin{theorem}[Alternating Series Test]
	Let $(a_n)$ be a sequence such that
	\begin{enumerate}
		\item $a_n \geq 0$ for all $n$,
		\item $a_{n+1} \leq a_n$ for all $n$ (i.e., $(a_n)$ is non-increasing), and
		\item $\lim_{n \to \infty} a_n = 0$.
	\end{enumerate}
	Then the alternating series
	\[
	\sum_{n=1}^{\infty} (-1)^{n+1} a_n
	\]
	converges.
\end{theorem}

\clearpage

\begin{heading}
	Problems with Conditional Convergence or Divergent Series
\end{heading}

\clearpage

\begin{heading}
	Rearrangements of Series
\end{heading}

\vspace{10pt}

\begin{definition}[Rearrangement of a Series]
	A \textit{rearrangement} of the series $\sum_{n=1}^{\infty} a_n$ is a series of the form
	\[
	\sum_{n=1}^{\infty} a_{p(n)},
	\]
	where $p: \mathbb{N} \to \mathbb{N}$ is a bijection.
\end{definition}


\clearpage

\begin{theorem}[Absolute Convergence and Rearrangements]
	Every rearrangement of an absolutely convergent series converges to the same sum as the original series.
\end{theorem}

\clearpage

\begin{theorem}[Riemann Rearrangements Theorem]
	If a series $\sum_{n=1}^{\infty} a_n$ is conditionally convergent, then for any real number $L$, there exists a rearrangement of the series that converges to $L$. Furthermore, there exists a rearrangement that diverges.	
\end{theorem}

\clearpage

\begin{heading}
	Double Summations and Products
\end{heading}

\vspace{10pt}

\begin{definition}[Double Series]
	A \textit{double series} is a series of the form
	\[
	\sum_{m=1}^{\infty} \sum_{n=1}^{\infty} a_{m,n},
	\]
	where $a_{m,n}$ are real numbers indexed by two indices $m$ and $n$.
\end{definition}

\clearpage

\begin{problem}
Consider the doubly indexed family of real numbers $(a_{ij})_{i,j \in \mathbb{N}}$
defined by
\[
a_{ij} =
\begin{cases}
-1, &\text{if } j=i,\\[6pt]
\dfrac{1}{2^{\,j-i}}, &\text{if } j>i,\\[6pt]
0, &\text{if } j<i.
\end{cases}
\]
Equivalently, the infinite matrix $(a_{ij})$ has the form
\[
\begin{pmatrix}
-1 & \tfrac{1}{2} & \tfrac{1}{4} & \tfrac{1}{8} & \tfrac{1}{16} & \cdots \\
0  & -1 & \tfrac{1}{2} & \tfrac{1}{4} & \tfrac{1}{8} & \cdots \\
0  & 0  & -1 & \tfrac{1}{2} & \tfrac{1}{4} & \cdots \\
0  & 0  & 0  & -1 & \tfrac{1}{2} & \cdots \\
\vdots & \vdots & \vdots & \vdots & \ddots & \ddots
\end{pmatrix}.
\]

We wish to attach a mathematical meaning to the formal double series
\[
\sum_{i=1}^{\infty} \sum_{j=1}^{\infty} a_{ij}.
\]
\end{problem}


\clearpage

\begin{theorem}[Fubini's Theorem for Series]
	If the double series $\sum_{m=1}^{\infty} \sum_{n=1}^{\infty} |a_{m,n}|$ converges, then the iterated sums
	\[
	\sum_{m=1}^{\infty} \left( \sum_{n=1}^{\infty} a_{m,n} \right) \quad \text{and} \quad \sum_{n=1}^{\infty} \left( \sum_{m=1}^{\infty} a_{m,n} \right)
	\]
	converge and are equal to the same value as the double series.
\end{theorem}

\clearpage

\begin{definition}[Product of Series]
	Given two series $\sum_{n=1}^{\infty} a_n$ and $\sum_{n=1}^{\infty} b_n$, their \textit{Cauchy product} is defined as the series
	\[
	\sum_{n=1}^{\infty} c_n,
	\]
	where
	\[
	c_n = \sum_{k=1}^{n} a_k b_{n-k+1}.
	\]
\end{definition}

\clearpage

\begin{theorem}[Mertens' Theorem]
	If $\sum_{n=1}^{\infty} a_n$ converges to $A$ and $\sum_{n=1}^{\infty} b_n$ converges to $B$, and at least one of these series converges absolutely, then their Cauchy product converges to $AB$.
\end{theorem}

\clearpage







\EndPage
\end{document}